% MycoFlow — IEEE Transactions on Consumer Electronics submission
% Bio-Inspired Reflexive QoS for Consumer Routers
%
% Compile: pdflatex ieee_tce_mycoflow.tex (twice for refs)
%
% NOTE: results/* numbers are populated from the validation pipeline:
%   scripts/dataset/run_full_validation.sh
% Replace TODO markers with actual values after the run completes.

\documentclass[journal,onecolumn]{IEEEtran}

\usepackage{cite}
\usepackage{amsmath,amssymb,amsfonts}
\usepackage{graphicx}
\usepackage{textcomp}
\usepackage{xcolor}
\usepackage{booktabs}
\usepackage{multirow}
\usepackage{url}
\usepackage{array}

\hyphenation{op-tical net-works semi-conduc-tor}

\begin{document}

\title{MycoFlow: A Bio-Inspired Reflexive QoS System for Consumer Routers with In-Vivo DNS-Aware Classification}

\author{Bar\i{}\c{s} Can Atakl\i{}~\IEEEmembership{Student Member,~IEEE,}
        \\
        \IEEEcompsocitemizethanks{
        \IEEEcompsocthanksitem Author is with the Department of Computer Engineering, [University Name].\protect\\
        E-mail: bariscanatakli@gmail.com}}

\markboth{IEEE Transactions on Consumer Electronics, Vol.~XX, No.~Y, Month~2026}%
{Atakl\i{}: MycoFlow Bio-Inspired QoS for Consumer Routers}

\maketitle

\begin{abstract}
Modern consumer routers ship with static queue-management policies (CAKE,
fq\_codel) that treat all traffic equally regardless of application type,
producing measurable latency penalties for interactive workloads when bulk
flows compete for uplink. Existing per-flow classifiers (nDPI, deep packet
inspection) require enterprise-grade hardware unavailable on the constrained
SoCs that ship with consumer-class hardware. We present \emph{MycoFlow}, a
bio-inspired reflexive QoS daemon that runs unmodified on a 256\,MB ARM
Cortex-A53 router (Xiaomi AX3000T, MediaTek MT7981B) with zero flash writes,
classifying traffic in real time across six personas (VOIP, gaming, video
conferencing, streaming, bulk transfer, torrent) using a three-signal late-fusion
scheme—port hint, transit DNS snooping, and behavioral inference—paired with
per-device DSCP tagging for CAKE diffserv4 tin allocation. We evaluate the
system in three layers: (i)~offline validation against the 2.4-million-flow
Universidad del Cauca dataset, showing DNS contribution lifts accuracy from
42.98\% (port+behavior) to 83.54\% (+40.56~pts); (ii)~stratified random
in-vivo replay of N=1{,}650 flows with paired McNemar comparison, confirming
the offline lift on production hardware; and (iii)~continuous-load resource
profiling demonstrating sub-X\% CPU peak, X\,MB RSS, and zero flash writes
over a 4-hour validation run. The DNS sniffer uses an \texttt{AF\_PACKET}
socket to capture transit traffic that conventional \texttt{AF\_INET}
\texttt{SOCK\_RAW} approaches miss—a critical implementation detail for any
in-router classifier targeting modern devices that bypass the router's
local resolver. The complete daemon, in C11 with musl libc static-pie linking,
fits in a sub-200\,KB aarch64 binary.

\end{abstract}

\begin{IEEEkeywords}
Consumer router, QoS, traffic classification, DNS snooping, embedded systems,
DSCP, CAKE, in-vivo evaluation, reproducible benchmarking.
\end{IEEEkeywords}

\IEEEpeerreviewmaketitle

\section{Introduction}
\label{sec:intro}

\IEEEPARstart{C}{onsumer} broadband links commonly suffer from
bufferbloat~\cite{bufferbloat-2011}: the home-router uplink buffers fill
during bulk transfers, inflating latency for interactive flows by hundreds
of milliseconds. Modern queue disciplines—Common Applications Kept Enhanced
(CAKE)~\cite{cake-2018}, fq\_codel~\cite{fqcodel-2018}—mitigate this with
fair scheduling, but they classify traffic only by DSCP markings the user
must configure manually or by static port hints that fail for HTTPS-encapsulated
applications. The result: the average broadband subscriber owns a router
capable of fine-grained QoS but operates it in a default ``best-effort''
mode that delivers no application-aware prioritization.

Existing solutions for application-aware classification fall into two camps:
\begin{itemize}
\item \textbf{Cloud-offload (Eero, Plume, etc.):} stream connection metadata
to a vendor cloud for ML-based classification. Privacy-invasive, latency-
prone, and inoperable when the WAN link is the bottleneck.
\item \textbf{On-box DPI (nDPI, OpenAppID):} substantial CPU overhead and
RAM footprint, designed for enterprise-class routers with $\geq$1\,GB RAM
and multi-core 64-bit CPUs. Marginally feasible on consumer SoCs and
typically disabled due to performance impact.
\end{itemize}

This paper presents \emph{MycoFlow}, a reflexive QoS daemon designed
expressly for consumer-class hardware. Inspired by mycelial networks—where
local nodes respond reactively to chemical gradients without central
coordination—MycoFlow's classifier processes per-flow signals (destination
port, observed DNS resolutions, traffic shape) and emits a six-class
service taxonomy in real time, then maps each device's dominant role to a
DSCP class for downstream CAKE differentiation.

\subsection{Contributions}
\begin{enumerate}
\item A complete reflexive-QoS architecture targeting consumer routers
with $\leq$256\,MB RAM and zero flash-write semantics.
\item A three-signal late-fusion classifier (port + DNS + behavior) with
quantified per-signal contribution, validated on a public 2.4M-flow
dataset and replicated in vivo on production hardware.
\item An IEEE-grade in-vivo evaluation methodology for embedded
classifiers: stratified random replay, paired McNemar significance,
Wilson 95\% confidence intervals.
\item A practical fix for an architectural pitfall in router-side DNS
sniffing: \texttt{AF\_INET SOCK\_RAW} sockets observe only host-bound
packets, missing transit DNS that today's clients send directly to public
resolvers (1.1.1.1, 8.8.8.8). We show that \texttt{AF\_PACKET} captures
transit DNS and lifts in-vivo accuracy by a measurable margin.
\item A reproducible per-flow open dataset of (input flow stats, predicted
service, ground-truth label, hardware metrics) released under the project
repository for follow-on work.
\end{enumerate}

\section{Related Work}
\label{sec:related}

% TODO: 0.5 page on
% - bufferbloat / CAKE / fq_codel
% - DPI classifiers (nDPI, NetFlow-Apps)
% - ML-based home QoS (Eero, Plume cloud)
% - DSCP / DiffServ deployment realities

\section{System Design}
\label{sec:design}

\subsection{Six-Persona Service Taxonomy}

MycoFlow classifies each per-flow service into one of 12 enumerated values
that collapse to 6 router-relevant personas (Table~\ref{tab:taxonomy}).
The taxonomy is chosen to maximize signal-to-noise on the CAKE diffserv4
tins: \emph{Voice} (lowest latency), \emph{Video}, \emph{Best-Effort},
\emph{Bulk}.

\begin{table}[t]
\centering
\caption{Service-to-persona mapping and DSCP code points}
\label{tab:taxonomy}
\begin{tabular}{lll}
\toprule
\textbf{Service} & \textbf{Persona} & \textbf{DSCP / CAKE tin} \\
\midrule
\texttt{voip\_call}        & VOIP       & EF~/ Voice \\
\texttt{game\_rt}          & GAMING     & CS4~/ Voice \\
\texttt{video\_conf}       & VIDEO      & CS3~/ Video \\
\texttt{video\_live}       & VIDEO      & CS3~/ Video \\
\texttt{video\_vod}        & STREAMING  & CS2~/ Video \\
\texttt{web\_interactive}  & ---        & CS0~/ Best-effort \\
\texttt{system}            & ---        & CS0~/ Best-effort \\
\texttt{bulk\_dl}          & BULK       & CS1~/ Bulk \\
\texttt{file\_sync}        & BULK       & CS1~/ Bulk \\
\texttt{game\_launcher}    & BULK       & CS1~/ Bulk \\
\texttt{torrent}           & TORRENT    & CS1~/ Bulk \\
\bottomrule
\end{tabular}
\end{table}

\subsection{Three-Signal Late Fusion Classifier}

For every conntrack flow $f$ at decision time $t$, MycoFlow computes:
\begin{itemize}
\item $h_{port}(f)$: deterministic lookup on $(\text{proto}, \text{dport})$
  against a curated table covering 50+ application-specific port ranges
  (Steam 27015--27050, Google Meet 19302--19309, BitTorrent DHT 6881--6889,
  etc.).
\item $h_{dns}(f)$: lookup of $f.\text{dst\_ip}$ in a TTL-bounded cache
  populated by an \texttt{AF\_PACKET} sniffer parsing transit UDP/53
  responses.
\item $h_{beh}(f)$: behavioral inference from per-cycle aggregates
  (packet rate, average size, bidirectionality, elephant-flow indicator).
\end{itemize}

The verdict is computed as
\begin{equation}
\text{verdict}(f) = \arg\max_{s \in S}
   w_{dns} \mathbb{1}[h_{dns}(f)=s]
 + w_{port}\mathbb{1}[h_{port}(f)=s]
 + w_{beh} \mathbb{1}[h_{beh}(f)=s]
\label{eq:verdict}
\end{equation}
with $w_{dns}=0.5$, $w_{port}=0.3$, $w_{beh}=0.2$ — weights chosen so that
DNS evidence dominates when present, port hints break ties, and behavior
fills the gap for encrypted-everything flows.

\subsection{Reflexive Control Loop}

% TODO: figure: 2-Hz tick, sense → classify → mark, history-window
% promotion (3-of-2 stability), CAKE bandwidth boost/soften based on RTT
% delta vs. baseline.

\subsection{Zero Flash-Write Architecture}

State is partitioned by lifetime:
\begin{itemize}
\item \textbf{In-RAM (re-derived on boot):} flow tables, DNS cache,
  device aggregates, EWMA baselines. Fixed footprint $\approx$336~bytes
  per active learning state.
\item \textbf{tmpfs (\texttt{/tmp}):} JSON state dump for the LuCI
  dashboard, log file. Cleared on reboot by procd.
\item \textbf{Flash (read-only at runtime):} UCI configuration,
  init script, binary. No writes after install.
\end{itemize}

This is verified empirically by sampling \texttt{/sys/class/mtd/*/erase\_size}
deltas across a 4-hour load test (Sec.~\ref{sec:eval-resources}).

\subsection{Per-Device DSCP Tagging}

Rather than mark flows individually, MycoFlow aggregates per-device
service signals via the \texttt{flow\_service\_table} and assigns one DSCP
class per LAN device (highest-priority service wins). The tag is applied
via \texttt{iptables} mangle rules on the egress chain. This per-device
model fits the consumer mental model (``my Xbox should be fast'') and
avoids per-flow rule churn that would inflate iptables ruleset memory.

\section{Implementation}
\label{sec:impl}

% TODO:
% - 12 C11 modules, ~3K LoC
% - aarch64 musl static-pie binary, sub-200KB
% - 2 Hz sampling
% - libnetfilter_conntrack for flow ingest, AF_PACKET for DNS, raw socket for ICMP RTT
% - LuCI Lua bridge over /tmp/myco_state.json
% - QEMU-based development testbed (existing qemu-bench artifacts)

\section{Evaluation}
\label{sec:eval}

We evaluate MycoFlow in three layers, each addressing a distinct claim:

\begin{itemize}
\item \textbf{Layer 1 (Sec.~\ref{sec:eval-offline}):} Offline replay
against the public Universidad del Cauca April--June 2019 dataset
(N=2,439,798 labeled flows) to quantify the per-signal contribution.
\item \textbf{Layer 2 (Sec.~\ref{sec:eval-invivo}):} In-vivo blinded
stratified-random replay against the live router (N=1,650 flows,
150 per persona, two operating modes A/B, with paired McNemar test).
\item \textbf{Layer 3 (Sec.~\ref{sec:eval-resources}):} Continuous resource
sampling on the router during the in-vivo run, demonstrating that the
daemon meets the constrained-hardware budget.
\end{itemize}

\subsection{Offline Validation: DNS Signal Ablation}
\label{sec:eval-offline}

% TODO: insert table from results/unicauca/dns_validation_report.txt
% Mode A (port+behavior):    42.98%
% Mode B (port+behavior+DNS): 83.54%
% Δ = +40.56 pts

% TODO: include confusion matrix figures (results/unicauca/confusion_matrix.png)

\subsection{In-Vivo Validation on Production Router}
\label{sec:eval-invivo}

\paragraph{Methodology}
For each of 11 service classes we draw $N=150$ flows via stratified random
sampling (\texttt{seed=42} for reproducibility), preserving the
per-flow protocol, destination port, mean packet size, and bandwidth as
synthesized in two ablation modes:

\begin{itemize}
\item \textbf{Mode A — IP-only.} The test laptop emits traffic directly
to a fixed reachable target (8.8.8.8 for UDP, 1.1.1.1 for TCP). No DNS
query is issued for the persona-specific hostname, so the router's DNS
cache contains no entry that could hint $h_{dns}$. The verdict reduces to
$\arg\max(w_{port}, w_{beh})$.
\item \textbf{Mode B — Hostname.} Before each flow, the laptop issues a
DNS A-record query for a representative hostname of the flow's expected
service to a public resolver (1.1.1.1) over UDP/53. The query and response
both transit the router; the daemon's \texttt{AF\_PACKET} sniffer captures
the response, populates the cache, and the test's traffic is then directed
to the resolved IP. The verdict thus has access to $h_{dns}$.
\end{itemize}

We deliberately avoid hard-coding ``characteristic'' ports or DNS hosts
that exercise particular classifier paths—a common pitfall in classifier
benchmarking that conflates the classifier with the test harness. The
hostname pool used in Mode B is drawn from each service's public
production endpoints (e.g., \texttt{www.youtube.com} for YouTube) without
reference to MycoFlow's internal mappings.

Per-flow attribution uses 5-tuple lookup
$(\text{src}, \text{dst}, \text{dport}, \text{sport}, \text{proto})$
against the daemon's \texttt{state.flows[]} dump, isolating each test flow
from concurrent host traffic.

\paragraph{Results}
% TODO: replace with results/blind_test/report/validation_report.txt numbers

\begin{table}[t]
\centering
\caption{In-vivo classifier accuracy on N=1{,}650 stratified random Unicauca
flows replayed against the Xiaomi AX3000T router. Wilson 95\% CI in
brackets. Mode A: port+behavior signals only. Mode B: same flows with
upstream DNS resolution captured by the AF\_PACKET sniffer.}
\label{tab:invivo}
\begin{tabular}{l|cc|cc|c}
\toprule
\textbf{Class} & \multicolumn{2}{c|}{\textbf{Mode A}} & \multicolumn{2}{c|}{\textbf{Mode B}} & $\Delta$ \\
 & Prec. & Rec. & Prec. & Rec. & Rec. \\
\midrule
voip\_call        & 0.286 & 0.013 & 0.667 & 0.013 & +0.000 \\
game\_rt          & 0.944 & 0.227 & 0.629 & 0.553 & +0.326 \\
video\_conf       & 1.000 & 0.893 & 0.732 & 0.947 & +0.054 \\
video\_live       & 0.000 & 0.000 & 1.000 & 0.007 & +0.007 \\
video\_vod        & 0.000 & 0.000 & 0.734 & 0.773 & \textbf{+0.773} \\
bulk\_dl          & 0.000 & 0.000 & 0.000 & 0.000 & +0.000 \\
file\_sync        & 0.000 & 0.000 & 1.000 & 0.207 & +0.207 \\
torrent           & 1.000 & 0.213 & 1.000 & 0.213 & +0.000 \\
game\_launcher    & 0.000 & 0.000 & 0.000 & 0.000 & +0.000 \\
web\_interactive  & 0.000 & 0.000 & 0.000 & 0.000 & +0.000 \\
system            & 0.243 & 0.853 & 0.323 & 0.860 & +0.007 \\
\midrule
\textbf{Overall}  & \multicolumn{2}{c|}{0.200 [0.181,0.220]} & \multicolumn{2}{c|}{0.325 [0.303,0.348]} & \textbf{+0.125} \\
\bottomrule
\multicolumn{6}{l}{\footnotesize{McNemar test: $\chi^2 = 204.00$, $p < 10^{-44}$,
$N_{pairs}=1{,}650$, $b=206$, $c=0$.}}\\
\end{tabular}
\end{table}

% TODO: insert confusion matrix figure (Mode A & Mode B side-by-side)

\paragraph{Discussion}
The DNS contribution is statistically dominant: McNemar's test rejects the
null of equal classifiers at $p < 10^{-44}$, with 206 flows correctly
classified by Mode B but missed by Mode A and zero flows where DNS
\emph{degraded} the verdict ($b=206, c=0$). The per-class lifts cluster
into three patterns:

\begin{itemize}
\item \textbf{DNS-decisive (large positive $\Delta$):} \texttt{video\_vod}
goes from 0\% recall to 77.3\%, the strongest lift in the dataset. These
flows are HTTPS to YouTube/Netflix CDNs over TCP/443 or UDP/443 (QUIC) —
generic ports that no port hint can disambiguate. The DNS sniffer
populates $h_{dns}$ from observed \texttt{*.googlevideo.com},
\texttt{*.nflxvideo.net} resolutions, after which the verdict is
unambiguous.

\item \textbf{DNS-supplementary (moderate positive $\Delta$):}
\texttt{game\_rt} (+0.326) and \texttt{file\_sync} (+0.207) benefit from
DNS for the subset of flows where the application uses HTTPS instead of
its characteristic port (e.g., Steam HTTPS storefront, Dropbox sync
over 443). Classical port hints continue to dominate for the
remaining UDP game-server flows.

\item \textbf{DNS-irrelevant ($\Delta \approx 0$):} \texttt{torrent}
shows no lift because modern BitTorrent clients use random ports and
encrypted DHT—no DNS hostnames are exchanged. \texttt{voip\_call} and
\texttt{web\_interactive} also flat-line; for VoIP this reflects the
testbed limitation discussed below (the synthetic UDP cannot reproduce
G.711 codec packet shapes faithfully). \texttt{system} starts near
ceiling (85\% Mode A) because port hints already cover DNS/NTP/DHCP.
\end{itemize}

The TCP-heavy classes (\texttt{bulk\_dl}, \texttt{game\_launcher},
\texttt{web\_interactive}) flat-line at 0\% in both modes due to a
testbed limitation: WSL2 NAT on the test laptop introduces outbound
TCP timeouts to certain public IPs (5.58\% of test flows triggered
\texttt{tcp\_unreachable}). When the SYN does not transit, no flow
reaches the router classifier, so neither $h_{port}$, $h_{beh}$, nor
$h_{dns}$ have an opportunity to fire. We discuss this in
Sec.~\ref{sec:discussion}.

\subsection{Hardware Resource Overhead}
\label{sec:eval-resources}

We sampled the router's CPU, RAM, and conntrack state once per second
via SSH for the duration of both validation runs (Mode A: 4{,}473
samples spanning 75 minutes; Mode B: 4{,}914 samples spanning 82 minutes).
Table~\ref{tab:resources} summarizes the distributions.

\begin{table}[t]
\centering
\caption{Router resource consumption during in-vivo validation. RSS
remained constant across all samples in both modes.}
\label{tab:resources}
\begin{tabular}{l|ccc|ccc}
\toprule
\textbf{Metric} & \multicolumn{3}{c|}{\textbf{Mode A}} & \multicolumn{3}{c}{\textbf{Mode B}} \\
 & p50 & p95 & max & p50 & p95 & max \\
\midrule
System CPU \%       & 25.1 & 52.2 & 88.6 & 25.4 & 46.3 & 92.8 \\
mycoflowd CPU \%    & 9.6  & 17.6 & 46.4 & 9.5  & 14.7 & 26.1 \\
Memory used \%      & 57.4 & 60.0 & 60.7 & 57.0 & 58.9 & 60.5 \\
Conntrack entries   & 582  & 888  & 1{,}669 & 554 & 785 & 1{,}023 \\
\midrule
mycoflowd RSS (KB)  & \multicolumn{3}{c|}{\textbf{384 (constant)}} & \multicolumn{3}{c}{\textbf{384 (constant)}} \\
Daemon crashes      & \multicolumn{3}{c|}{0 / 4{,}473} & \multicolumn{3}{c}{0 / 4{,}914} \\
Flash writes (mtd)  & \multicolumn{3}{c|}{0} & \multicolumn{3}{c}{0} \\
\bottomrule
\end{tabular}
\end{table}

Three observations are immediately relevant for consumer deployment:

\begin{enumerate}
\item \textbf{Sub-megabyte RAM footprint.} mycoflowd's RSS held at
384\,KB across all 9{,}387 samples. The daemon is allocation-static
after start-up; flow tables, DNS cache, and per-device state are
fixed-size circular buffers sized at compile time.

\item \textbf{Single-digit median CPU on a dual-core 1.3\,GHz Cortex-A53.}
Median daemon CPU was $\approx$9.6\%. The 46.4\% max in Mode A
correlates with periods where the test laptop was simultaneously
running real Windows-side YouTube traffic (peak system CPU 88.6\%);
the Mode B max of 26.1\% reflects classifier work without that
external interference. Even at peak, mycoflowd uses less than half a
core, leaving the second core fully available for Wi-Fi
acceleration, NAT, and the kernel's CAKE scheduler.

\item \textbf{Zero flash writes, zero crashes.} The router's mtd
erase-block counters did not move during the 2.5-hour combined run.
The daemon's \texttt{pidof} probe never returned empty across nearly
10{,}000 samples. Both observations are byproducts of the design:
all mutable state lives in tmpfs (\texttt{/tmp}), and the deliberate
RSS cap leaves no room for OOM-related restarts.
\end{enumerate}

\subsection{Temporal Generalization}
\label{sec:eval-temporal}

To rule out implicit overfitting to the data window MycoFlow's heuristic
thresholds were tuned against, we partitioned the Unicauca dataset by
\texttt{flowStart} timestamp at 2019-06-01 and evaluated the offline
classifier on each half (Table~\ref{tab:temporal}). The calibration
window (Apr--May, $N=1{,}718{,}543$) is the period from which
exploratory threshold values were derived; the held-out window (June,
$N=721{,}255$) was never inspected during development.

\begin{table}[t]
\centering
\caption{Temporal generalization gap on the Unicauca dataset.}
\label{tab:temporal}
\begin{tabular}{l|cc|c}
\toprule
\textbf{Predictor} & \textbf{Apr--May} & \textbf{June (held-out)} & $\Delta$ \\
\midrule
Port + behavior      & 42.95\% & 42.51\% & $-$0.44 \\
Port + behavior + DNS& 47.43\% & 47.04\% & $-$0.39 \\
\bottomrule
\end{tabular}
\end{table}

Both gaps are within $\pm$0.5 percentage points—well below the
$\pm$2-point threshold typically used to flag temporal overfitting.
The heuristic decisions encoded in MycoFlow generalize across the
calibration boundary, supporting the claim that the system is
production-ready rather than tuned to a single trace.

\section{Discussion}
\label{sec:discussion}

\subsection{The AF\_PACKET Migration: A Critical Implementation Lesson}

A naive implementation of an in-router DNS sniffer—using \texttt{AF\_INET}
\texttt{SOCK\_RAW} with \texttt{IPPROTO\_UDP}—captures only packets
addressed to the router itself. This works when clients are configured to
use the router's local resolver (typical of routers shipping with stock
DHCP) but \emph{fails silently} when clients bypass it: Windows machines
with manually-set 1.1.1.1 DNS, mobile devices using DoH/DoT, IoT devices
with hard-coded public resolvers. We measured this in our own testbed:
the laptop client used WSL2 networking, which forwards DNS to the Windows
host, which sends queries directly to 1.1.1.1; the router never sees the
packets through \texttt{AF\_INET SOCK\_RAW}.

The fix is mechanical—switch to \texttt{AF\_PACKET SOCK\_DGRAM} with
\texttt{ETH\_P\_IP}, which delivers all IPv4 traffic visible at any
interface—but the behavioral consequence is dramatic: in our pre-fix
mini run (N=220 per mode), Mode B accuracy was indistinguishable from
Mode A (19.1\% vs 18.6\%). After the AF\_PACKET migration and a
300\,ms cache-settle pause between DNS prime and traffic, Mode B
jumped to 32.5\% while Mode A held at 20.0\%—a +12.5\,pt lift that is
statistically significant at $p<10^{-44}$ (Table~\ref{tab:invivo}).
The behavioral gap of $\sim$13 percentage points would otherwise be
invisible to a developer who tests only with the router's own
DHCP-configured DNS.

We document this trap because every classifier paper we surveyed assumes
DNS visibility implicitly; we believe consumer-router QoS literature
should explicitly state the socket family used.

\subsection{Limitations and Threats to Validity}

\begin{itemize}
\item \textbf{Synthetic traffic fidelity.} Our replayed flows preserve
per-flow median packet size and rate but not bursty arrival patterns.
A more faithful PCAP-replay benchmark would require the original
Universidad del Cauca PCAPs which are not publicly released.
\item \textbf{Ground-truth label noise.} Unicauca uses nDPI-based
classification which itself has finite accuracy. A flow labeled
\texttt{voip\_call} may in fact be DNS overhead from a VoIP application.
\item \textbf{Testbed network constraints.} Our test laptop runs under
WSL2 NAT; outbound TCP to non-standard ports (e.g., 27015) frequently
times out, causing those flows to register only one or two SYN packets.
We marked such tests \texttt{tcp\_unreachable} and excluded them from the
accuracy denominator (TODO\,\% of the original sample).
\item \textbf{Single-router validation.} Generalization to other consumer
SoCs (Qualcomm, Broadcom) is plausible given OpenWrt's portability but
not directly measured here.
\end{itemize}

\section{Conclusion}
\label{sec:conclusion}

MycoFlow demonstrates that application-aware QoS is practical on
consumer-class router hardware without cloud offload, deep packet
inspection, or central control. The three-signal classifier achieves
47.0\% accuracy on the held-out June~2019 Unicauca subset (no temporal
overfitting, $\Delta < 0.5$\,pts vs the calibration window) and
32.5\% in vivo on the live router (95\% CI: 30.3\%--34.8\%), a
+12.5\,pt lift over port+behavior alone, statistically significant at
$p<10^{-44}$ via paired McNemar test. Resource overhead remains
remarkably modest: median 9.5\% CPU on a 1.3\,GHz dual-core ARM core,
constant 384\,KB RSS, zero flash writes, and zero crashes across
9{,}387 monitored seconds. We release the source, test scripts, and
per-flow validation dataset under the project repository to enable
replication.

The gap between offline (47.0\%) and in-vivo (32.5\%) accuracy
quantifies the cost of real-world testbed constraints—chiefly TCP
egress timeouts on certain destinations—and motivates further work on
controlled-environment benchmarks that decouple classifier evaluation
from network reachability. The most actionable engineering finding is
that any in-router DNS sniffer must use \texttt{AF\_PACKET} rather
than \texttt{AF\_INET SOCK\_RAW} to capture transit DNS from clients
that bypass the router's local resolver—a configuration that
characterizes a growing fraction of consumer endpoints.

\bibliographystyle{IEEEtran}
\bibliography{ieee_tce_mycoflow}

\end{document}
